% GENERATED FILE. Edit canonical lessons, then run npm run generate:books.
% canonical-source-hash: fnv1a64:5372ad000d6f66c9
% canonical-lessons: JA-C31-dasu, JA-C31-otosu, JA-C31-naosu, JA-C31-ugoku, JA-C31-gohan

\chapter{To take out, To drop, To fix, To move, Rice}
\label{ch:ja-to-take-out-to-drop-to-fix-to-move-rice}
\hlchaptermodality{31}

\begin{chapteropening}
\textbf{By the end of this chapter:} \emph{I can say to take out, to drop, to fix, to move and rice.}

Complete the last lesson of chapter 31: i can say to take out, to drop, to fix, to move and rice.
\end{chapteropening}

\section[dasu]{\ja{だす} (dasu) — to take out}

\label{lesson:JA-C31-dasu}



\begin{quote}
Before the new one: say the Japanese for to lend, then the Japanese for to give back.
\end{quote}

\subsection*{You'll want to know: \ja{だす}}
\textbf{\ja{だす}} — \emph{dasu} — \textquotedblleft{}to take out\textquotedblright{}.

\begin{grammarlens}[title={What you've built}]
One more word for this chapter.
\end{grammarlens}

\subsection*{Your turn}
\begin{itemize}
\raggedright
  \item \emph{Say it:} \emph{dasu}
  \item \emph{Say it:} \emph{dasu}, once more
  \item \emph{Say it:} say \emph{dasu}
  \item \emph{Recall:} say the Japanese for to lend, then the Japanese for to give back, then say \emph{dasu} again
\end{itemize}

\begin{tcolorbox}[breakable,colback=teal!4,colframe=teal!35!black,title={Before you move on}]
What does \ja{だす} mean? (\textquotedblleft{}To take out\textquotedblright{}.) Say it once more.
\end{tcolorbox}
\section[otosu]{\ja{おとす} (otosu) — to drop}

\label{lesson:JA-C31-otosu}



\begin{quote}
Before the new one: say the Japanese for to give back, then the Japanese for to take out.
\end{quote}

\subsection*{You'll want to know: \ja{おとす}}
\textbf{\ja{おとす}} — \emph{otosu} — \textquotedblleft{}to drop\textquotedblright{}.

\begin{grammarlens}[title={What you've built}]
One more word for this chapter.
\end{grammarlens}

\subsection*{Your turn}
\begin{itemize}
\raggedright
  \item \emph{Say it:} \emph{otosu}
  \item \emph{Say it:} \emph{otosu}, once more
  \item \emph{Say it:} say \emph{otosu}
  \item \emph{Recall:} say the Japanese for to give back, then the Japanese for to take out, then say \emph{otosu} again
\end{itemize}

\begin{tcolorbox}[breakable,colback=teal!4,colframe=teal!35!black,title={Before you move on}]
What does \ja{おとす} mean? (\textquotedblleft{}To drop\textquotedblright{}.) Say it once more.
\end{tcolorbox}
\section[naosu]{\ja{なおす} (naosu) — to fix}

\label{lesson:JA-C31-naosu}



\begin{quote}
Before the new one: say the Japanese for to take out, then the Japanese for to drop.
\end{quote}

\subsection*{You'll want to know: \ja{なおす}}
\textbf{\ja{なおす}} — \emph{naosu} — \textquotedblleft{}to fix\textquotedblright{}.

\begin{grammarlens}[title={What you've built}]
One more word for this chapter.
\end{grammarlens}

\subsection*{Your turn}
\begin{itemize}
\raggedright
  \item \emph{Say it:} \emph{naosu}
  \item \emph{Say it:} \emph{naosu}, once more
  \item \emph{Say it:} say \emph{naosu}
  \item \emph{Recall:} say the Japanese for to take out, then the Japanese for to drop, then say \emph{naosu} again
\end{itemize}

\begin{tcolorbox}[breakable,colback=teal!4,colframe=teal!35!black,title={Before you move on}]
What does \ja{なおす} mean? (\textquotedblleft{}To fix\textquotedblright{}.) Say it once more.
\end{tcolorbox}
\section[ugoku]{\ja{うごく} (ugoku) — to move}

\label{lesson:JA-C31-ugoku}



\begin{quote}
Before the new one: say the Japanese for to drop, then the Japanese for to fix.
\end{quote}

\subsection*{You'll want to know: \ja{うごく}}
\textbf{\ja{うごく}} — \emph{ugoku} — \textquotedblleft{}to move\textquotedblright{}.

\begin{grammarlens}[title={What you've built}]
One more word for this chapter.
\end{grammarlens}

\subsection*{Your turn}
\begin{itemize}
\raggedright
  \item \emph{Say it:} \emph{ugoku}
  \item \emph{Say it:} \emph{ugoku}, once more
  \item \emph{Say it:} say \emph{ugoku}
  \item \emph{Recall:} say the Japanese for to drop, then the Japanese for to fix, then say \emph{ugoku} again
\end{itemize}

\begin{tcolorbox}[breakable,colback=teal!4,colframe=teal!35!black,title={Before you move on}]
What does \ja{うごく} mean? (\textquotedblleft{}To move\textquotedblright{}.) Say it once more.
\end{tcolorbox}
\section[gohan]{\ja{ごはん} (gohan) — rice; a meal}

\label{lesson:JA-C31-gohan}



\begin{quote}
Before the new one: say the Japanese for to fix, then the Japanese for to move.
\end{quote}

\subsection*{You'll want to know: \ja{ごはん}}
\textbf{\ja{ごはん}} — \emph{gohan} — \textquotedblleft{}rice; a meal\textquotedblright{}.

\begin{grammarlens}[title={What you've built}]
One more word for this chapter.
\end{grammarlens}

\subsection*{Your turn}
\begin{itemize}
\raggedright
  \item \emph{Say it:} \emph{gohan}
  \item \emph{Say it:} \emph{gohan}, once more
  \item \emph{Say it:} say \emph{gohan}
  \item \emph{Recall:} say the Japanese for to fix, then the Japanese for to move, then say \emph{gohan} again
\end{itemize}

\begin{tcolorbox}[breakable,colback=teal!4,colframe=teal!35!black,title={Before you move on}]
What does \ja{ごはん} mean? (\textquotedblleft{}Rice; a meal\textquotedblright{}.) Say it once more.
\end{tcolorbox}
